\documentclass[12pt]{article}
\usepackage[utf8]{inputenc}
\usepackage[T1]{fontenc}
\usepackage{hyperref}
\title{Wunambal Language Resources for K-12 Curriculum Development}
\author{Compiled by Claude Sonnet 5.5 (claude-sonnet-5-5) for Andr\'as Kornai}
\date{}
\begin{document}
\maketitle

\section{Introduction}
Wunambal (ISO 639-3 wub, Glottocode wuna1249) is a Worrorran language of the north-west Kimberley region of Western Australia. It is very closely related to Gaambera. AIATSIS describes the two as differing very little in pronunciation and vocabulary. Wikipedia reports fewer than ten speakers in 2009.

\section{Primary Educational Resources}
The search found a federal grant, ``Strengthening Wunambal Gaambera Language and Culture'' (about AUD 220,000, running to June 2025). It is reported to fund digital and physical learning resources, including dictionaries, language lessons, videos and multimedia storybooks. The grant record is at \url{https://www.grants.gov.au/Ga/Show/7422f780-f244-42a1-91bd-0360edd03563}. No finished classroom materials were seen in the search results.

\section{Linguistic and Reference Materials}
\begin{itemize}
\item Therese Carr, \emph{Wunambal: A language of the North-West Kimberley Region, Western Australia} (2001 thesis, University of New England): \url{https://rune.une.edu.au/web/handle/1959.11/16787}
\item Wunambal language (Wikipedia): \url{https://en.wikipedia.org/wiki/Wunambal_language}
\item AUSTLANG record for Gaambera (K39): \url{https://collection.aiatsis.gov.au/austlang/language/K39}
\item AUSTLANG entry for the \emph{Handbook of Kimberley Languages} (1988): \url{https://collection.aiatsis.gov.au/austlang/language/k22}
\end{itemize}
Secondary sources mention \emph{Uunguu plants and animals} (Karadada et al., 2011), a dictionary of biological knowledge that treats Wunambal and Gaambera together. They also mention H.H.J. Coate's 1948 manuscript \emph{English-Wunambal dictionary}. No URL for either was seen, and the manuscript was not found online.

\section{Religious and Community Resources}
No scripture was found. The grant above refers to a Wunambal Gaambera community language effort, but no community organisation page was seen in the results.

\section{Digital and Audio Resources}
No online audio, video or dictionary was found, although the grant is meant to produce such material.

\section{Curriculum Development Context}
Few speakers remain. The Carr thesis and the Uunguu dictionary are the best starting points. Any school curriculum would need to be built together with the Wunambal Gaambera community and based on the materials the revival grant produces.
\end{document}